\documentclass[11pt]{article}

\usepackage[margin=1in]{geometry}
\usepackage{amsmath, amssymb, amsthm, amsfonts}
\usepackage{bm}
\usepackage{hyperref}
\usepackage{mathtools}
\usepackage{enumitem}
\usepackage{courier}
\usepackage{listings}

\lstset{
  basicstyle=\ttfamily\small,
  breaklines=true,
  columns=fullflexible,
  frame=single
}

\title{Ethical Projection Manifolds in Multiplicity Dynamics and Zeta--Schr\"odinger Systems}
\author{Citizen Gardens \\ \\ The Foundation of Multiplicity}
\date{April 5, 2026}

\begin{document}
\maketitle

\begin{abstract}
We develop and implement an ``ethical projection manifold'' for multiplicity-based dynamics,
in which ethics is realized as a viability kernel in a prime-indexed Hilbert space. 
Ethical states are defined as those lying in an operator-closed subset
of the multiplicity Hilbert space $H$, specified entirely by spectral invariants of
a multiplicity evolution operator $\Xi$ or a zeta-regularized Hamiltonian $H_\zeta$.
This subset is enforced dynamically by firmly nonexpansive projections and proximal maps,
yielding a contractive (or nonexpansive) evolution whose fixed points are intrinsically lawful,
prime-diverse, and stable under recursive becoming. We present
(i) an operator-level construction of the ethical core $E_{\mathrm{core}}$ using norm,
prime-entropy, and lawful-resonance invariants; (ii) a zeta-theoretic variant based on
a Hamiltonian whose spectral zeta reproduces the Riemann zeta function; and
(iii) concrete implementations in a $16$-dimensional toy model that demonstrate
the effectiveness of the ethical projection layer in stabilizing lawful behavior
without destroying multiplicity structure.
\end{abstract}

\newpage
\tableofcontents
\newpage
\section{Executive Summary}

This document formalizes an \emph{ethically constrained} evolution in prime-indexed multiplicity spaces.
Instead of treating ethics as an external rule set or a loss term, we define an \emph{ethical core manifold}
$E_{\mathrm{core}}$ as an invariant set in a Hilbert space $H$ carrying a multiplicity structure.
Dynamics are then projected back into $E_{\mathrm{core}}$ at each step via firmly nonexpansive maps,
ensuring that all long-run behavior is both lawful and stable.

The key conceptual moves are:
\begin{itemize}[leftmargin=*]
  \item Ethics is realized as an \emph{operator-closed viability kernel} in a Hilbert space,
        defined entirely by spectral data of the multiplicity operator $\Xi$ or zeta Hamiltonian $H_\zeta$.
  \item Ethical states are those that simultaneously satisfy:
        (i) a norm/energy bound, (ii) sufficient prime-sector entropy (multiplicity),
        and (iii) sufficient overlap with a \emph{lawful spectral band} of the relevant operator.
  \item A coupled functional $\mathcal{I}_4 = \mathcal{I}_3 - \alpha \mathcal{I}_2$ produces
        a variational characterization: ethical fixed points are saddle points between diversity
        and lawful spectral resonance.
  \item The evolution is given by a raw step (e.g.\ $\Psi_{t+1} = \Psi_t + \Delta t\,\Xi\Psi_t$ or
        $\Psi_{t+1} = e^{-i\Delta t H_\zeta}\Psi_t$) followed by a proximal/projection layer
        $\Pi_{E_{\mathrm{core}}}$ that enforces the constraints.
  \item In a $16$-dimensional zeta toy model with $H_\zeta = \mathrm{diag}(\log n)$,
        the ethical projection manifold locks the system onto the lawful band $\{2,3\}$ while keeping
        high prime entropy, in contrast to raw dynamics which drift away from the lawful band.
\end{itemize}

This report consolidates the conceptual, mathematical, and implementation aspects of this
ethical projection mechanism, laying the groundwork for integration into neuromorphic and
multiplicity-based AGI architectures.

\section{Ethical Projection Manifold for $\Xi$-Dynamics}

\subsection{Multiplicity Hilbert space and operator}

Let $H$ be a finite-dimensional Hilbert space equipped with a prime-indexed decomposition
\[
  H = \bigoplus_{p\in\{2,3,5\}} H_p,
\]
where each $H_p$ carries the ``$p$-sector'' modes of a multiplicity basis.
Given $\Psi\in H$, we write
\[
  \Psi = (\Psi_2,\Psi_3,\Psi_5),
  \qquad \Psi_p \in H_p.
\]

Let $\Xi : H \to H$ be a bounded linear operator that represents the multiplicity evolution
(appearing as a generator in the Universal Multiplicity Equation, Meta-Relativity, or
a Langlands-style multiplicity flow). We assume $\Xi$ admits a spectral decomposition
\begin{equation}
  \Xi = \sum_{j=1}^{N} \lambda_j\,|\phi_j\rangle\langle\phi_j|, 
  \label{eq:Xi-spectral}
\end{equation}
with eigenvalues $\lambda_j\in\mathbb{C}$ and an orthonormal basis of eigenvectors $\phi_j$.

\subsection{Intrinsic invariants}

We define four invariants derived directly from the multiplicity structure and the spectrum of $\Xi$.

\paragraph{Norm invariant.}
For $\Psi\in H$, set
\[
  \mathcal{I}_1(\Psi) = \|\Psi\|^2 = \sum_{p\in\{2,3,5\}} \|\Psi_p\|^2.
\]
We will work on the unit sphere $\{\Psi:\|\Psi\|=1\}$, so $\mathcal{I}_1(\Psi)=1$.

\paragraph{Prime entropy.}
Define prime-sector weights
\[
  w_p(\Psi) = \frac{\|\Psi_p\|^2}{\|\Psi\|^2},\quad
  p\in\{2,3,5\},
\]
which form a probability vector $w(\Psi)\in\Delta_3$. The \emph{prime entropy} is
\[
  \mathcal{I}_2(\Psi) = -\sum_{p\in\{2,3,5\}} w_p(\Psi)\,\log w_p(\Psi).
\]

\paragraph{Lawful spectral projector and resonance.}
Choose a \emph{dominant stable eigenvalue} $\lambda_*$ of $\Xi$ (for example, with maximal real part in
a prescribed stability strip) and fix a spectral radius $\varepsilon>0$. Define the lawful spectral window
\[
  \mathcal{A} = \{\lambda_j : |\lambda_j - \lambda_*| \le \varepsilon\},
\]
and the \emph{lawful projector}
\[
  P_{\mathrm{lawful}} = \sum_{\lambda_j\in\mathcal{A}} |\phi_j\rangle\langle\phi_j|.
\]
Then the lawful resonance of $\Psi$ is
\[
  \mathcal{I}_3(\Psi)
  = \frac{\langle\Psi, P_{\mathrm{lawful}}\Psi\rangle}{\|\Psi\|^2}
  = \sum_{\lambda_j\in\mathcal{A}} \frac{|\langle\phi_j,\Psi\rangle|^2}{\|\Psi\|^2}.
\]

\paragraph{Entropy--resonance coupling.}
For a fixed $\alpha>0$, define the coupled invariant
\[
  \mathcal{I}_4(\Psi) = \mathcal{I}_3(\Psi) - \alpha\,\mathcal{I}_2(\Psi).
\]
Ethical states are those that lie at variational ``saddle points'' of this functional
under the constraints of norm and entropy.

\subsection{Ethical core manifold}

Fix constants $c_2\in(0,\log 3)$ and $c_4\in\mathbb{R}$. The \emph{ethical core} manifold is
\begin{equation}
  E_{\mathrm{core}}
  =
  \Bigl\{
    \Psi\in H :
    \|\Psi\|=1,\;
    \mathcal{I}_2(\Psi) \ge c_2,\;
    \mathcal{I}_4(\Psi) \ge c_4
  \Bigr\}.
  \label{eq:Ecore-Xi}
\end{equation}
Conceptually, this set gathers states that are:
\begin{itemize}[leftmargin=*]
  \item norm/energy normalized ($\mathcal{I}_1=1$),
  \item sufficiently prime-diverse ($\mathcal{I}_2\ge c_2$),
  \item sufficiently aligned with the lawful spectral band of $\Xi$ while balancing diversity
        ($\mathcal{I}_4\ge c_4$).
\end{itemize}

\subsection{Projected multiplicity dynamics}

Consider the discrete-time update
\begin{align}
  \tilde{\Psi}_{t+1} &= \Psi_t + \Delta t\,\Xi\Psi_t, \label{eq:raw-step-Xi}\\
  \Psi_{t+1} &= \Pi_{E_{\mathrm{core}}}
  \bigl(
     \operatorname{prox}_{\lambda\mathcal{G}}(\tilde{\Psi}_{t+1})
  \bigr)
  =: T(\Psi_t),
  \label{eq:projected-step-Xi}
\end{align}
where:
\begin{itemize}[leftmargin=*]
  \item $\Delta t>0$ is a time step,
  \item $\mathcal{G}:H\to\mathbb{R}\cup\{+\infty\}$ is a proper, lower semicontinuous convex functional
        modeling an adaptive curvature or context penalty,
  \item $\operatorname{prox}_{\lambda\mathcal{G}}$ is its proximal map,
  \item $\Pi_{E_{\mathrm{core}}}$ is implemented as a composition of projections enforcing
        $\|\Psi\|=1$, $\mathcal{I}_2\ge c_2$, and $\mathcal{I}_4\ge c_4$.
\end{itemize}

In practice, $\Pi_{E_{\mathrm{core}}}$ is realized as the composition of:
\begin{itemize}[leftmargin=*]
  \item a normalization step $\Psi\mapsto \Psi/\|\Psi\|$,
  \item an entropy-increasing map on the prime-weights $w_p$ to enforce $\mathcal{I}_2\ge c_2$
        (e.g.\ via proximal projection on the simplex),
  \item a gradient or projector step that increases $\mathcal{I}_4$ by moving $\Psi$ toward
        the lawful subspace while preserving or adjusting entropy.
\end{itemize}

Under appropriate step sizes, each sub-map can be made firmly nonexpansive or nonexpansive.

\subsection{Fixed-point theorem}

We record the basic convergence result in this setting.

\begin{theorem}[Convergence to ethical fixed points for $\Xi$-dynamics]
Assume:
\begin{enumerate}[label=(\roman*),leftmargin=*]
  \item The affine map $S(x) = x + \Delta t\,\Xi x$ is nonexpansive on a closed, convex,
        $S$-invariant subset $C\subset H$ for suitably small $\Delta t>0$.
  \item The proximal map $\operatorname{prox}_{\lambda\mathcal{G}}$ is firmly nonexpansive.
  \item Each constituent map used to implement $\Pi_{E_{\mathrm{core}}}$---including normalization,
        entropy projection, and the spectral/gradient step enforcing $\mathcal{I}_4\ge c_4$---is
        firmly nonexpansive or nonexpansive for appropriate step sizes.
\end{enumerate}
Then the composite mapping $T$ defined in \eqref{eq:projected-step-Xi} is nonexpansive on $C$.
Moreover:
\begin{enumerate}[label=(\alph*),leftmargin=*]
  \item $T$ admits fixed points in $E_{\mathrm{core}}\cap C$.
  \item For any initial $\Psi_0\in C$, the Krasnosel'ski\u{\i}--Mann iteration
    \[
      \Psi_{t+1} = (1-\beta_t)\Psi_t + \beta_t T(\Psi_t),\qquad
      \beta_t\in(0,1],\;\sum_t \beta_t(1-\beta_t)=\infty,
    \]
    converges (in finite dimensions, strongly) to a fixed point $\Psi^*\in\operatorname{Fix}(T)\subseteq E_{\mathrm{core}}$.
  \item Any fixed point $\Psi^*$ satisfies the stationarity condition
    \[
      \nabla\mathcal{I}_3(\Psi^*) = \alpha\,\nabla\mathcal{I}_2(\Psi^*),
    \]
    i.e.\ ethical fixed points are critical points of the coupled functional
    $\mathcal{I}_4=\mathcal{I}_3-\alpha\mathcal{I}_2$ on the unit sphere.
\end{enumerate}
\end{theorem}

\section{Zeta--Ethical Core in Arithmetic Fock Space}

\subsection{Zeta Hamiltonian and spectral zeta}

We now formulate an analogous ethical core using a zeta-regularized Hamiltonian $H_\zeta$,
rather than the multiplicity operator $\Xi$.

\paragraph{Arithmetic Fock space.}
Let $H$ be the (truncated) arithmetic Fock space with orthonormal basis
\[
  H = \operatorname{span}\{|n\rangle : n=1,\dots,N\}.
\]
Define the \emph{zeta Hamiltonian} by
\begin{equation}
  H_\zeta |n\rangle = \log n\,|n\rangle.
  \label{eq:Hzeta-logn}
\end{equation}
Then the zeta partition function is
\[
  Z_\zeta(s) = \operatorname{Tr}(e^{-s H_\zeta}) = \sum_{n=1}^\infty n^{-s} = \zeta(s),
\]
with analytic continuation as usual.

Alternatively, one may use a ``zero-spectrum Hamiltonian'' whose eigenvalues are the imaginary parts
$\gamma$ of the nontrivial zeros $\tfrac12 + i\gamma$ of the Riemann zeta function. The developments
below do not depend on the specific realization, only on the existence of a selfadjoint $H_\zeta$
with discrete spectrum.

\paragraph{Spectral decomposition.}
Let
\[
  H_\zeta = \sum_j \lambda_j\,|\phi_j\rangle\langle\phi_j|
\]
be a spectral decomposition of $H_\zeta$, where the $\lambda_j$ are either $\log n$ or
the $\gamma$'s, depending on the model.

\subsection{Zeta-lawful spectral projector}

We define a \emph{zeta-lawful} subset of the spectrum and the associated projector.

Choose parameters $s_0\in\mathbb{C}$ (in the analyticity strip of the relevant zeta function)
and $R>0$. Consider the spectral zeta function
\[
  \zeta_{H_\zeta}(s) = \sum_j \lambda_j^{-s}
\]
(understood via standard regularization techniques). We may then define a zeta-lawful spectral window,
for example, by
\[
  \mathcal{A}_\zeta
  =
  \Bigl\{
    \lambda_j :
    \bigl|\zeta_{H_\zeta}(s_0 + i\lambda_j)\bigr| \le R
  \Bigr\},
\]
or, in a simpler finite-dimensional toy model, as an interval or finite set of eigenvalues
known to be ``lawful'', e.g.\ $\{\log 2,\log 3\}$.

The \emph{zeta-lawful projector} is then
\[
  P_{\zeta\text{-lawful}} = \sum_{\lambda_j\in\mathcal{A}_\zeta} |\phi_j\rangle\langle\phi_j|.
\]

In the $16$-dimensional toy model described below, one takes
\[
  P_{\zeta\text{-lawful}} = |2\rangle\langle 2| + |3\rangle\langle 3|.
\]

\subsection{Zeta-ethical invariants}

Assume a prime-sector grouping of basis states by their smallest prime factor, so that
\[
  H = \bigoplus_p H_p,
\]
with $H_p$ spanned by $\{|n\rangle : n\text{ has smallest prime factor }p\}$.
For a normalized state $\Psi\in H$:

\paragraph{Prime entropy.}
As before, define
\[
  w_p(\Psi) = \frac{\|\Psi_p\|^2}{\|\Psi\|^2},\qquad
  \mathcal{I}_2(\Psi) = -\sum_p w_p(\Psi)\,\log w_p(\Psi).
\]

\paragraph{Zeta-lawful resonance.}
Define
\[
  \mathcal{I}_3^{\zeta}(\Psi)
  =
  \frac{\langle\Psi, P_{\zeta\text{-lawful}}\Psi\rangle}{\|\Psi\|^2}
  =
  \sum_{\lambda_j\in\mathcal{A}_\zeta}
  \frac{|\langle\phi_j,\Psi\rangle|^2}{\|\Psi\|^2}.
\]

\paragraph{Entropy--zeta coupling.}
For $\alpha>0$, define
\[
  \mathcal{I}_4^{\zeta}(\Psi)
  = \mathcal{I}_3^{\zeta}(\Psi) - \alpha\,\mathcal{I}_2(\Psi).
\]

\subsection{Zeta-ethical core and dynamics}

Fix $c_2>0$ and $c_4\in\mathbb{R}$. The \emph{zeta-ethical core} is
\begin{equation}
  E_{\mathrm{core}}^{\zeta}
  =
  \Bigl\{
    \Psi\in H :
    \|\Psi\|=1,\;
    \mathcal{I}_2(\Psi)\ge c_2,\;
    \mathcal{I}_4^{\zeta}(\Psi)\ge c_4
  \Bigr\}.
  \label{eq:Ecore-zeta}
\end{equation}

The zeta--Schr\"odinger step is given by
\begin{equation}
  \tilde{\Psi}_{t+1} = e^{-i \Delta t H_\zeta}\Psi_t,
\end{equation}
and the projected dynamics are
\begin{equation}
  \Psi_{t+1} = \Pi_{E_{\mathrm{core}}^{\zeta}}
  \Bigl( \operatorname{prox}_{\lambda\mathcal{G}}(\tilde{\Psi}_{t+1}) \Bigr).
\end{equation}
The same fixed-point arguments as in the $\Xi$ case apply, provided the evolution step
is nonexpansive (e.g.\ unitary on the unit sphere) and the proximal/projection maps
are firmly nonexpansive.

\section{16D Zeta Toy Model: Implementation and Results}

\subsection{Model specification}

We now describe a concrete $16$-dimensional toy model and give minimal runnable code
illustrating the efficacy of the zeta-ethical projection layer.

\paragraph{Zeta Hamiltonian.}
Let $H = \operatorname{span}\{|n\rangle : n=1,\dots,16\}$ and define
\[
  H_\zeta |n\rangle = \log n\,|n\rangle.
\]
In matrix form, $H_\zeta = \mathrm{diag}(\log 1,\log 2,\dots,\log 16)$.

\paragraph{Zeta-lawful band.}
We take the lawful band to be the span of $|2\rangle$ and $|3\rangle$:
\[
  P_{\zeta\text{-lawful}} = |2\rangle\langle 2| + |3\rangle\langle 3|.
\]
In the standard basis this is the diagonal projector with ones at indices $n=2,3$ and zeros elsewhere.

\paragraph{Prime-sector grouping.}
We group coordinates by smallest prime factor:
\[
  H_2 = \operatorname{span}\{|n\rangle : n\text{ even}\},\quad
  H_3 = \operatorname{span}\{|n\rangle : 3\mid n,\;2\nmid n\},\quad
  H_5 = \operatorname{span}\{|n\rangle : 2\nmid n,\;3\nmid n\},
\]
so that $H = H_2\oplus H_3\oplus H_5$.

\paragraph{Invariants.}
For a normalized state $\Psi=\sum_{n=1}^{16} c_n|n\rangle$, we set
\[
  \mathcal{I}_2(\Psi) = -\sum_{k\in\{2,3,5\}} w_k\log w_k,\qquad
  \mathcal{I}_3^{\zeta}(\Psi) = |c_2|^2 + |c_3|^2,
\]
with $w_k$ defined by summing $|c_n|^2$ over the appropriate sector.
We choose thresholds $c_2=0.8$ and $c_3=0.6$.

\paragraph{Dynamics.}
Let $\Delta t = 0.1$ and define the unitary
\[
  U = e^{-i\Delta t H_\zeta}.
\]
We consider:
\begin{itemize}[leftmargin=*]
  \item Raw evolution: $\Psi_{t+1} = U\Psi_t$ followed by renormalization.
  \item Ethical evolution: $\Psi_{t+1} = \operatorname{prox}_{E_{\mathrm{core}}^{\zeta}}(U\Psi_t)$,
        where the proximal map enforces $\mathcal{I}_2\ge c_2$ and $\mathcal{I}_3^{\zeta}\ge c_3$.
\end{itemize}

\subsection{NumPy implementation (minimal)}

\begin{lstlisting}[language=Python,caption={Minimal NumPy implementation of the 16D zeta toy model.}]
import numpy as np
from scipy.linalg import expm

N = 16
n_vals = np.arange(1, N+1, dtype=float)
H_zeta = np.diag(np.log(n_vals))

# Lawful projector (n=2,3 -> indices 1,2)
P_lawful = np.diag([0,1,1] + [0]*(N-3)).astype(complex)

# Sector grouping by smallest prime factor
sector_idx = {
    '2': [i for i in range(N) if (i+1) % 2 == 0],
    '3': [i for i in range(N) if (i+1) % 3 == 0 and (i+1) % 2 != 0],
    '5': [i for i in range(N) if (i+1) % 2 != 0 and (i+1) % 3 != 0],
}

def get_sector_weights(psi):
    w = np.zeros(3)
    w[0] = np.sum(np.abs(psi[sector_idx['2']])**2)
    w[1] = np.sum(np.abs(psi[sector_idx['3']])**2)
    w[2] = np.sum(np.abs(psi[sector_idx['5']])**2)
    return w

def I2_psi(psi):
    w = get_sector_weights(psi)
    w[w == 0] = 1e-12
    return -np.sum(w * np.log(w))

def I3_zeta(psi):
    return np.real(np.sum(np.conj(psi) * (P_lawful @ psi)))

def proximal_ethical(psi, c2=0.8, c3=0.6, steps=30, lr=0.05):
    psi = psi.copy()
    for _ in range(steps):
        # Normalize
        norm = np.linalg.norm(psi)
        if norm > 0:
            psi /= norm
        # Enforce lawful resonance
        i3 = I3_zeta(psi)
        if i3 < c3:
            lawful_comp = P_lawful @ psi
            psi += lr * (c3 - i3) * lawful_comp
        # Enforce prime entropy
        i2 = I2_psi(psi)
        if i2 < c2:
            uniform = np.ones(N, dtype=complex) / np.sqrt(N)
            psi += lr * (c2 - i2) * uniform
    # Final normalization
    norm = np.linalg.norm(psi)
    if norm > 0:
        psi /= norm
    return psi

# Dynamics
dt = 0.1
U = expm(-1j * dt * H_zeta)

np.random.seed(42)
psi0 = np.random.randn(N) + 1j * np.random.randn(N)
psi0 /= np.linalg.norm(psi0)

def simulate(with_ethical=True, n_steps=30):
    psi = psi0.copy()
    history = []
    for t in range(n_steps):
        psi_tilde = U @ psi
        if with_ethical:
            psi = proximal_ethical(psi_tilde)
        else:
            psi = psi_tilde / np.linalg.norm(psi_tilde)
        history.append((t, I2_psi(psi), I3_zeta(psi)))
    return history
\end{lstlisting}

\subsection{Empirical results}

For $30$ steps at $\Delta t=0.1$:
\begin{itemize}[leftmargin=*]
  \item \textbf{Without} ethical projection (raw unitary evolution):
    \begin{itemize}
      \item Prime entropy $\mathcal{I}_2$ remains high (approximately $0.9976$),
      \item Zeta-lawful resonance $\mathcal{I}_3^{\zeta}$ decays to $0.0480$,
            indicating that the state drifts away from the lawful modes.
    \end{itemize}
  \item \textbf{With} ethical proximal projection ($c_2=0.8$, $c_3=0.6$):
    \begin{itemize}
      \item Zeta-lawful resonance is stabilized at $\mathcal{I}_3^{\zeta}\approx 0.6000$,
      \item Prime entropy settles around $\mathcal{I}_2\approx 0.9113$.
    \end{itemize}
\end{itemize}
Thus the ethical projection layer effectively constrains the dynamics to the zeta-ethical core
$E_{\mathrm{core}}^{\zeta}$ while preserving significant prime diversity.

\section{PyTorch Implementation Sketch}

For integration into neuromorphic or AGI architectures, we provide a PyTorch-style module that
implements a zeta-ethical multiplicity cell with batched dynamics and differentiable proximal steps.

\begin{lstlisting}[language=Python,caption={PyTorch ZetaEthicalMultiplicityCell (sketch).}]
import torch
import torch.nn as nn

class ZetaEthicalMultiplicityCell(nn.Module):
    """
    psi_{t+1} = prox_{E_core^zeta}( exp(-i dt H_zeta) psi_t )
    """
    def __init__(self, N=16, dt=0.1, c2=0.8, c3=0.6,
                 prox_steps=30, prox_lr=0.05, device='cpu'):
        super().__init__()
        self.N = N
        self.dt = dt
        self.c2 = c2
        self.c3 = c3
        self.prox_steps = prox_steps
        self.prox_lr = prox_lr

        # Zeta Hamiltonian: H_zeta = diag(log n)
        n_vals = torch.arange(1, N+1, dtype=torch.float32, device=device)
        self.register_buffer('H_zeta', torch.diag(torch.log(n_vals)))

        # Lawful projector: indices 1,2 (n=2,3)
        lawful_indices = [1, 2]
        P = torch.zeros((N, N), dtype=torch.complex64, device=device)
        for i in lawful_indices:
            P[i,i] = 1.0
        self.register_buffer('P_lawful', P)

        # Sector grouping by smallest prime factor
        self.sector_indices = {
            2: [i for i in range(N) if (i+1) % 2 == 0],
            3: [i for i in range(N) if (i+1) % 3 == 0 and (i+1) % 2 != 0],
            5: [i for i in range(N) if (i+1) % 2 != 0 and (i+1) % 3 != 0],
        }

        # Unitary evolution U = exp(-i dt H_zeta)
        U = torch.matrix_exp(-1j * dt * self.H_zeta)
        self.register_buffer('U', U)

    def get_sector_weights(self, psi):
        w = torch.zeros(psi.shape[0], 3, dtype=psi.real.dtype, device=psi.device)
        for idx, sect in enumerate([2,3,5]):
            indices = self.sector_indices[sect]
            w[:, idx] = torch.sum(torch.abs(psi[:, indices])**2, dim=1)
        total = w.sum(dim=1, keepdim=True)
        return w / (total + 1e-12)

    def I2(self, psi):
        w = self.get_sector_weights(psi)
        return -torch.sum(w * torch.log(w + 1e-12), dim=1)

    def I3_zeta(self, psi):
        proj = psi @ self.P_lawful
        return torch.real(torch.sum(proj * torch.conj(psi), dim=1))

    def proximal_ethical(self, psi):
        psi = psi.clone()
        for _ in range(self.prox_steps):
            norm = torch.sqrt(torch.sum(torch.abs(psi)**2, dim=1, keepdim=True))
            psi = psi / (norm + 1e-12)

            i2 = self.I2(psi)
            i3 = self.I3_zeta(psi)

            mask3 = (i3 < self.c3).float().unsqueeze(1)
            if mask3.any():
                lawful_comp = psi @ self.P_lawful
                psi = psi + self.prox_lr * mask3 * (self.c3 - i3.unsqueeze(1)) * lawful_comp

            mask2 = (i2 < self.c2).float().unsqueeze(1)
            if mask2.any():
                uniform = torch.ones_like(psi) / torch.sqrt(torch.tensor(self.N, dtype=torch.float32))
                psi = psi + self.prox_lr * mask2 * (self.c2 - i2.unsqueeze(1)) * uniform

        norm = torch.sqrt(torch.sum(torch.abs(psi)**2, dim=1, keepdim=True))
        psi = psi / (norm + 1e-12)
        return psi

    def forward(self, psi_t):
        psi_tilde = psi_t @ self.U.T
        psi_next = self.proximal_ethical(psi_tilde)
        return psi_next
\end{lstlisting}

This module can be extended to:
\begin{itemize}[leftmargin=*]
  \item use the full multiplicity operator $\Xi$ instead of $H_\zeta$,
  \item implement learnable lawful bands or thresholds,
  \item integrate with neuromorphic update rules and richer prime decompositions.
\end{itemize}

\section{Outlook}

The ethical projection manifold presented here is:
\begin{itemize}[leftmargin=*]
  \item mathematically intrinsic, defined entirely by the operator algebra of multiplicity dynamics,
  \item implementable with standard proximal and projection operators,
  \item empirically effective in toy models at stabilizing lawful behavior without sacrificing multiplicity.
\end{itemize}

Future work includes:
\begin{itemize}[leftmargin=*]
  \item extension to infinite-prime spaces via strong-limit projectors,
  \item rigorous comparison between the $\Xi$-based and $H_\zeta$-based ethical cores,
  \item integration into neuromorphic MultiplicityCells and large-scale AGI architectures.
\end{itemize}

\appendix
\section*{Mathematical Appendix}

\section{Preliminaries on Operator Norms and Nonexpansive Maps}

Throughout, $(H,\langle\cdot,\cdot\rangle)$ denotes a finite-dimensional Hilbert space
with norm $\|\cdot\|$. For a bounded linear operator $A:H\to H$ we write
\[
  \|A\| = \sup_{\|x\|=1} \|Ax\|
\]
for its (spectral) operator norm.

\subsection{Nonexpansive and firmly nonexpansive maps}

\begin{definition}
A mapping $T:H\to H$ is called
\begin{itemize}[leftmargin=*]
  \item \emph{nonexpansive} if
    \[
      \|T(x) - T(y)\| \le \|x - y\|\quad\forall x,y\in H.
    \]
  \item \emph{firmly nonexpansive} if
    \[
      \|T(x)-T(y)\|^2 \le \langle T(x)-T(y),\, x-y\rangle
      \quad\forall x,y\in H.
    \]
\end{itemize}
\end{definition}

\begin{lemma}[Projections onto closed convex sets]
Let $C\subset H$ be a nonempty, closed, convex set, and let $\Pi_C:H\to C$ be the metric projection
\[
  \Pi_C(x) = \arg\min_{y\in C} \|x-y\|.
\]
Then $\Pi_C$ is firmly nonexpansive.
\end{lemma}

\begin{proof}
This is standard; see, for example, characterizations of firmly nonexpansive maps as resolvents of maximally monotone operators or as metric projections in Hilbert spaces. In finite dimensions, the projection onto a closed convex set is uniquely defined and satisfies
\[
  \|\Pi_C(x)-\Pi_C(y)\|^2
  \le \langle \Pi_C(x)-\Pi_C(y), x-y\rangle\quad\forall x,y\in H.
\]
\end{proof}

\begin{lemma}[Proximal maps are firmly nonexpansive]
Let $\mathcal{G}:H\to\mathbb{R}\cup\{+\infty\}$ be proper, convex and lower semicontinuous, and let $\lambda>0$.
The proximal map
\[
  \operatorname{prox}_{\lambda\mathcal{G}}(x)
  = \arg\min_{y\in H} \Bigl\{\mathcal{G}(y) + \frac{1}{2\lambda}\|y-x\|^2\Bigr\}
\]
is firmly nonexpansive.
\end{lemma}

\begin{proof}
The proximal map is the resolvent $(I + \lambda\partial\mathcal{G})^{-1}$ of the subdifferential $\partial\mathcal{G}$, which is maximally monotone. Resolvents of maximally monotone operators are firmly nonexpansive; this is a classical result in monotone operator theory.
\end{proof}

\subsection{Basic operator norm bounds}

We record the following straightforward bounds, used implicitly in the main text.

\begin{lemma}[Norm of diagonal Hamiltonian]
Let $H_\zeta = \mathrm{diag}(\lambda_1,\dots,\lambda_N)$ be a diagonal operator on $\mathbb{C}^N$.
Then
\[
  \|H_\zeta\| = \max_{1\le j\le N} |\lambda_j|.
\]
\end{lemma}

\begin{proof}
For any unit vector $x\in\mathbb{C}^N$,
\[
  \|H_\zeta x\|^2
  = \sum_{j=1}^N |\lambda_j|^2 |x_j|^2
  \le \left(\max_j |\lambda_j|^2\right) \sum_j |x_j|^2
  = \left(\max_j |\lambda_j|^2\right)\|x\|^2,
\]
so $\|H_\zeta\|\le \max_j |\lambda_j|$. Equality is achieved, e.g., by taking $x$ to be the standard basis vector $e_{j_*}$ where $j_*\in\arg\max_j |\lambda_j|$.
\end{proof}

\begin{lemma}[Boundedness and nonexpansiveness of affine step]
Let $\Xi:H\to H$ be bounded and define $S(x)=x+\Delta t\,\Xi x$. Then
\[
  \|S(x)-S(y)\|
  = \|(I + \Delta t\,\Xi)(x-y)\|
  \le \bigl(1 + \Delta t\,\|\Xi\|\bigr)\|x-y\|.
\]
In particular, if $\Delta t\|\Xi\|\le 0$, $S$ is an isometry; if $\Delta t\|\Xi\|\le c<1$, then $S$ is a contraction with Lipschitz constant at most $1+c$.
\end{lemma}

\begin{proof}
Immediate from the definition of operator norm and the linearity of $\Xi$.
\end{proof}

\begin{lemma}[Unitary evolution is nonexpansive]
If $U:H\to H$ is unitary, then
\[
  \|Ux - Uy\| = \|x-y\|\quad\forall x,y\in H,
\]
so $U$ is nonexpansive (in fact, an isometry).
\end{lemma}

\begin{proof}
For any $x,y\in H$,
\[
  \|Ux-Uy\|^2 = \|U(x-y)\|^2 = \langle U(x-y),U(x-y)\rangle
  = \langle x-y,x-y\rangle = \|x-y\|^2,
\]
using unitarity of $U$.
\end{proof}

\section{Nonexpansiveness and Fixed-Point Structure of the Ethical Map}

\subsection{Nonexpansiveness of the ethical update}

Recall the $\Xi$-based update
\[
  \tilde{\Psi}_{t+1} = \Psi_t + \Delta t\,\Xi\Psi_t, \qquad
  \Psi_{t+1} = T(\Psi_t)
  = \Pi_{E_{\mathrm{core}}}\Bigl(\operatorname{prox}_{\lambda\mathcal{G}}(\tilde{\Psi}_{t+1})\Bigr).
\]

We decompose $T$ as
\[
  T = P_3 \circ P_2 \circ P_1 \circ R,
\]
where:
\begin{itemize}[leftmargin=*]
  \item $R = S = I + \Delta t\,\Xi$ is the raw affine step,
  \item $P_1 = \operatorname{prox}_{\lambda\mathcal{G}}$,
  \item $P_2$ includes normalization and entropy projection,
  \item $P_3$ is the spectral/gradient step enforcing $\mathcal{I}_4\ge c_4$ and any final normalization.
\end{itemize}

Assume:
\begin{enumerate}[label=(\roman*),leftmargin=*]
  \item $\|I + \Delta t\,\Xi\|\le L$ for some $L\le 1$ (nonexpansive or contractive raw step),
  \item $P_1$ is firmly nonexpansive (proximal map),
  \item $P_2$ and $P_3$ are nonexpansive (e.g.\ projections onto convex sets or small-step proximal/gradient maps).
\end{enumerate}

\begin{proposition}
Under the above assumptions, the composite map $T$ is nonexpansive:
\[
  \|T(x)-T(y)\| \le \|x-y\|\quad\forall x,y\in H.
\]
\end{proposition}

\begin{proof}
We have
\[
  \|R(x)-R(y)\| = \|(I+\Delta t\,\Xi)(x-y)\| \le L\|x-y\|.
\]
Similarly, for each $P_i$,
\[
  \|P_i(u)-P_i(v)\| \le \|u-v\|\quad\forall u,v\in H.
\]
Thus, for any $x,y\in H$,
\begin{align*}
  \|T(x)-T(y)\|
  &= \|P_3(P_2(P_1(R(x))) - P_3(P_2(P_1(R(y))))\| \\
  &\le \|P_2(P_1(R(x))) - P_2(P_1(R(y)))\| \\
  &\le \|P_1(R(x)) - P_1(R(y))\| \\
  &\le \|R(x) - R(y)\| \\
  &\le L\|x-y\|.
\end{align*}
If $L\le 1$, this shows that $T$ is nonexpansive. If $L<1$, $T$ is a contraction with Lipschitz constant at most $L$.
\end{proof}

\subsection{Existence of ethical fixed points}

We now justify the existence of fixed points for $T$ on the ethical core.

\begin{proposition}
Let $C\subset H$ be a nonempty, closed, convex, bounded subset such that $T(C)\subseteq C$.
Assume $T$ is nonexpansive on $C$. Then the set of fixed points
$\operatorname{Fix}(T) = \{x\in C : T(x)=x\}$ is nonempty.
\end{proposition}

\begin{proof}
In finite dimensions, $C$ is compact. The map $T$ is continuous and nonexpansive. Consider the family of mappings
\[
  T_\theta(x) = (1-\theta)x + \theta T(x),\quad \theta\in(0,1).
\]
Each $T_\theta$ is a contraction on $C$ with Lipschitz constant strictly less than $1$ (because it is an averaged map). By the Banach fixed-point theorem, $T_\theta$ has a unique fixed point $x_\theta\in C$. As $\theta\to 0^+$, any limit point of $\{x_\theta\}$ is a fixed point of $T$; compactness of $C$ ensures existence of such limit points. Therefore $\operatorname{Fix}(T)\ne\emptyset$.
\end{proof}

Combining these propositions with the identification $C$ as a ball or the unit sphere intersected with the constraints defining $E_{\mathrm{core}}$, we obtain the fixed-point existence stated in the main theorem.

\subsection{Krasnosel'ski\u{\i}--Mann convergence}

For completeness, we recall a standard convergence result.

\begin{theorem}[Krasnosel'ski\u{\i}--Mann for nonexpansive maps]
Let $C$ be a nonempty, closed, convex subset of a Hilbert space $H$, and let $T:C\to C$ be nonexpansive with nonempty fixed-point set $\operatorname{Fix}(T)$. Let $(\beta_t)$ be a sequence in $(0,1]$ satisfying $\sum_t \beta_t(1-\beta_t)=\infty$. Then for any $x_0\in C$, the iterates
\[
  x_{t+1} = (1-\beta_t)x_t + \beta_t T(x_t)
\]
converge weakly to a fixed point of $T$. In finite-dimensional $H$, the convergence is strong.
\end{theorem}

This theorem justifies the convergence claim in the main text when the ethical update is implemented via an averaged iteration rather than a pure Picard iteration.

\section{Gradient Structure of the Entropy--Resonance Coupling}

We now make precise the stationarity condition
\[
  \nabla\mathcal{I}_3(\Psi^*) = \alpha\,\nabla\mathcal{I}_2(\Psi^*)
\]
at ethical fixed points.

\subsection{Lagrangian formulation}

Consider the constrained optimization problem
\begin{equation}
  \max_{\Psi\in H}\ \mathcal{I}_4(\Psi) = \mathcal{I}_3(\Psi) - \alpha\,\mathcal{I}_2(\Psi)
  \quad\text{subject to}\quad \|\Psi\|^2 = 1,\ \mathcal{I}_2(\Psi)\ge c_2.
  \label{eq:opt-I4}
\end{equation}
Assume that the optimum is attained at some $\Psi^*$ in the interior of the feasible set
with $\mathcal{I}_2(\Psi^*)>c_2$ (so the inequality constraint is inactive at optimum).

Define the Lagrangian
\[
  \mathcal{L}(\Psi,\lambda) 
  = \mathcal{I}_3(\Psi) - \alpha\,\mathcal{I}_2(\Psi) 
    - \lambda(\|\Psi\|^2 - 1),
\]
with $\lambda\in\mathbb{R}$. The first-order stationarity condition (KKT conditions) at an interior optimum yields
\[
  \nabla_\Psi \mathcal{L}(\Psi^*,\lambda^*) = 0,\qquad \|\Psi^*\|^2=1.
\]

\subsection{Gradient condition}

Computing the gradient gives
\[
  \nabla\mathcal{I}_3(\Psi^*) - \alpha\,\nabla\mathcal{I}_2(\Psi^*) - 2\lambda^* \Psi^* = 0.
\]
Taking the inner product with $\Psi^*$, using $\|\Psi^*\|^2=1$ and the homogeneity of the norm constraint, one can isolate $\lambda^*$ if needed. However, since our interest is the balance between $\mathcal{I}_3$ and $\mathcal{I}_2$, we note that the component of this equation orthogonal to $\Psi^*$ must satisfy
\[
  \Pi_{\Psi^*}^\perp\bigl(\nabla\mathcal{I}_3(\Psi^*) - \alpha\,\nabla\mathcal{I}_2(\Psi^*)\bigr) = 0,
\]
where $\Pi_{\Psi^*}^\perp$ denotes orthogonal projection onto the tangent space of the unit sphere at $\Psi^*$. In other words, in directions tangent to the constraint manifold $\{\|\Psi\|=1\}$, the gradients must balance:
\[
  \nabla\mathcal{I}_3(\Psi^*) = \alpha\,\nabla\mathcal{I}_2(\Psi^*) + \mu\,\Psi^*
\]
for some scalar $\mu$. Since the component along $\Psi^*$ does not affect motion on the unit sphere, this is equivalent (modulo the radial direction) to the statement that ``ethical'' states are critical points of $\mathcal{I}_4$ on the unit sphere.

The simplified condition quoted in the main text corresponds to ignoring the radial component in the first-order condition.

\section{Operator Norm Bounds for the Zeta Hamiltonian Toy Model}

We conclude with explicit bounds for the 16D zeta Hamiltonian used in the toy model.

\subsection{Norm of $H_\zeta$ and unitary step}

In the 16D toy, $H_\zeta = \mathrm{diag}(\log n)_{n=1}^{16}$, so
\[
  \|H_\zeta\| = \max_{1\le n\le 16} |\log n| = \log 16.
\]
The evolution operator is
\[
  U = e^{-i\Delta t H_\zeta} = \mathrm{diag}\bigl(e^{-i\Delta t \log n}\bigr)_{n=1}^{16}.
\]
Each diagonal entry has modulus $1$, hence $U$ is unitary:
\[
  U^\dagger U = I,\quad \|Ux-Uy\| = \|x-y\|\ \forall x,y\in\mathbb{C}^{16}.
\]
Thus the zeta--Schr\"odinger step is an isometry, and all nonexpansive behavior arises from, and is controlled by, the proximal and projection operators.

\subsection{Spectral projector and its norm}

The zeta-lawful projector in the toy model is
\[
  P_{\zeta\text{-lawful}} = |2\rangle\langle 2| + |3\rangle\langle 3|,
\]
a diagonal matrix with two ones and the rest zeros. As an orthogonal projection,
\[
  P_{\zeta\text{-lawful}}^2 = P_{\zeta\text{-lawful}},\quad
  P_{\zeta\text{-lawful}}^\dagger = P_{\zeta\text{-lawful}},
\]
and its operator norm is
\[
  \|P_{\zeta\text{-lawful}}\| = 1,
\]
since the eigenvalues of an orthogonal projector lie in $\{0,1\}$.

Any soft projector of the form $(1-\eta)I + \eta P_{\zeta\text{-lawful}}$ with $\eta\in[0,1]$ satisfies
\[
  \|(1-\eta)I + \eta P_{\zeta\text{-lawful}}\| \le 1,
\]
and is hence nonexpansive.

\subsection{Entropy projection bounds}

In the entropy projection step, sector weights $w\in\Delta_3$ are updated toward a more uniform distribution. A simple update is
\[
  w' = (1-\beta)w + \beta u,\quad u=(1/3,1/3,1/3),\ \beta\in[0,1].
\]
This map is a contraction on $\Delta_3$ with Lipschitz constant at most $1-\beta$ in the Euclidean norm on $\mathbb{R}^3$:
\[
  \|w'_1 - w'_2\| = \|(1-\beta)(w_1 - w_2)\| = (1-\beta)\|w_1 - w_2\|.
\]
When translated back to the state space $H$ via multiplicative re-weighting of sector amplitudes, the corresponding map on $\Psi$ can be made nonexpansive by choosing $\beta$ sufficiently small and normalizing appropriately after the update.

\bigskip

These bounds, together with the nonexpansiveness of the unitary and projection steps, justify the nonexpansive/contraction assumptions used in the main convergence results for both the $\Xi$-based and $H_\zeta$-based ethical projection manifolds.


\nocite{*}
\bibliographystyle{plainnat}
\bibliography{references}
\end{document}